\paragraph{Why hydraulic memory differs between soils.} Water released from storage leaves by the faster of two routes. In coarse soils (large $n$, high $K_s$) the conductivity at field capacity is high, so released water drains before roots and the surface can use it; in fine soils the conductivity is low and the same water stays near the roots and the surface, where it is transpired or evaporated.
\paragraph{Why the memory transfers to other experiments.} The same competition decides where any extra water goes, whether it comes from initial storage, from a lower evaporation factor or from a mulch; the share is a property of the soil's conductivity, not of the source of the water.
\paragraph{Why warming raises and wetter climates lower the recoverable share.} Warming raises the evaporative demand, so the withdrawal side of the competition is faster and less of a saved millimetre drains. A wetter climate keeps the profile closer to field capacity, where drainage is fast, so more of a saving drains.
\paragraph{Why nitrogen has no effect.} The nitrogen rate changes the crop's growth but not the hydraulic properties, and in these runs it changes the water fluxes too little to move the shares; it therefore serves as a negative control.
\paragraph{Why an evaporative term is unattributed.} The reported components cover transpiration, soil evaporation and canopy evaporation; evaporation of ponded water, litter and snow is part of evapotranspiration but not of those components, so it appears as a residual in the run table.
\paragraph{Why evapotranspiration follows a compact law and drainage does not.} Evapotranspiration is limited by energy and by water and is smooth in its inputs over 20 years. Drainage is a threshold process (water drains only above field capacity) whose annual total depends on the timing of rain and irrigation events, which the long-term means do not contain.
\paragraph{Why judging a saving against soil evaporation alone misleads.} A mulch changes soil evaporation and also the unattributed evaporation; if only the reported soil evaporation is used as the denominator, the drainage share of the change is computed against the wrong total and can change sign.

